\documentclass[11pt]{amsart}
\usepackage[letterpaper,margin=1in]{geometry}
\usepackage{lmodern}
\usepackage{amsmath,amssymb,amsthm,mathtools}
\usepackage{enumitem}
\usepackage{needspace}
\usepackage{microtype}
\usepackage[hidelinks]{hyperref}

\hypersetup{
  pdftitle={A divisor-closed criterion for Behrend sequences},
  pdfauthor={Will Howe, Grant Huffman, Jensen Kohlmeyer, Liam Kruer},
  pdfsubject={A necessary and sufficient condition for a set of multiples to have natural density one},
  pdfkeywords={Behrend sequences, sets of multiples, natural density, divisor-closed sets}
}

\newtheorem{theorem}{Theorem}[section]
\newtheorem{proposition}[theorem]{Proposition}
\newtheorem{lemma}[theorem]{Lemma}
\newtheorem{corollary}[theorem]{Corollary}
\theoremstyle{definition}
\newtheorem{definition}[theorem]{Definition}
\numberwithin{equation}{section}

\newcommand{\N}{\mathbb N}
\newcommand{\Nzero}{\mathbb N_0}
\newcommand{\Prob}{\mathbf P}
\newcommand{\Exp}{\mathbf E}
\newcommand{\one}{\mathbf 1}
\newcommand{\Hblock}[2]{\mathcal H_{#1,#2}}
\DeclareMathOperator{\lcm}{lcm}
\DeclareMathOperator{\Div}{Div}

\title[Divisor-closed families and Behrend sequences]{A divisor-closed criterion for Behrend sequences}
\author[Howe, Huffman, Kohlmeyer, and Kruer]{Will Howe, Grant Huffman, Jensen Kohlmeyer, and Liam Kruer}
\date{November 11, 2026}
\subjclass[2020]{11B05}
\keywords{Behrend sequences, sets of multiples, natural density, divisor-closed sets}
\thanks{\texttt{howew@purdue.edu} (W.~Howe), \texttt{huffma97@purdue.edu} (G.~Huffman),
\texttt{jenwin@purdue.edu} (J.~Kohlmeyer), and \texttt{kruerl@purdue.edu} (L.~Kruer).}

\begin{document}

\begin{abstract}
We give a necessary and sufficient condition for a set of positive integers to be a Behrend sequence.
The condition quantifies over sequences of finite divisor-closed sets supported on consecutive prime blocks.
For each element at one level, its permitted extensions at the next level retain a prescribed reciprocal mass.
We prove that the set of multiples of $A$ has natural density one if and only if $A$ meets every such family.
The proof combines the classical Davenport--Erd\H{o}s finite-approximation theorem with a finite pruning argument for independent geometric prime valuations.
\end{abstract}

\maketitle

\section{Statement of the criterion}

Write $\N=\{1,2,\ldots\}$ and $\Nzero=\{0,1,2,\ldots\}$.
For $A\subseteq\N$, let
\begin{equation}\label{eq:multiples}
  M_A=\{n\in\N: a\mid n\text{ for some }a\in A\}.
\end{equation}
For $B\subseteq\N$, its lower natural density is
\[
  \underline d(B)=\liminf_{x\to\infty}\frac{|B\cap[1,x]|}{x}.
\]
When the corresponding limit exists, it is denoted by $d(B)$.
A set $A$ is called a \emph{Behrend sequence} when $d(M_A)=1$.
Erd\H{o}s asked for a necessary and sufficient condition for this property~\cite[p.~77]{Erdos1979}.

Let $p_1<p_2<\cdots$ be the primes, and let $v_p(n)$ denote the exponent of $p$ in $n$.
For $\ell\in\Nzero$, define the prime-prefix projection
\begin{equation}\label{eq:projection}
  \pi_\ell(n)=\prod_{i\leq\ell}p_i^{v_{p_i}(n)}.
\end{equation}
For integers $0\leq r\leq s$, put
\begin{equation}\label{eq:blocks}
  \Hblock{r}{s}
  =\{n\in\N: v_{p_i}(n)=0\text{ whenever }i\notin(r,s]\},
  \qquad
  Q_{r,s}=\prod_{r<i\leq s}\left(1-\frac1{p_i}\right).
\end{equation}
Empty products equal one, so $\pi_0(n)=1$, $\Hblock{r}{r}=\{1\}$, and $Q_{r,r}=1$.
A set $S\subseteq\N$ is \emph{divisor closed} if $n\in S$ and $d\mid n$ imply $d\in S$.

\begin{theorem}\label{thm:main}
For every $A\subseteq\N$, the following are equivalent.
\begin{enumerate}[label=\textup{(\roman*)},leftmargin=2.4em,itemsep=3pt]
  \item The set $M_A$ has natural density one.
  \item For every admissible family $(\ell_k,S_k)_{k\geq0}$ in Definition~\ref{def:admissible},
  \begin{equation}\label{eq:criterion}
    A\cap\bigcup_{k\geq0}S_k\neq\varnothing.
  \end{equation}
\end{enumerate}
\end{theorem}

\Needspace{19\baselineskip}
\begin{definition}\label{def:admissible}
An \emph{admissible family} consists of integers
\[
  0\leq\ell_0<\ell_1<\ell_2<\cdots
\]
and finite sets $S_k\subseteq\N$, $k\in\Nzero$, satisfying $S_0=\{1\}$ and the following conditions for every $k\geq1$.
\begin{enumerate}[label=\textup{(\roman*)},leftmargin=2.4em,itemsep=3pt]
  \item\label{ad:support} $S_k$ is divisor closed and $S_k\subseteq\Hblock{\ell_0}{\ell_k}$.
  \item\label{ad:projection} $\pi_{\ell_{k-1}}(S_k)\subseteq S_{k-1}$.
  \item\label{ad:retention} For every $m\in S_{k-1}$, the finite set
  \begin{equation}\label{eq:fiber}
    F_k(m)=\{n\in\Hblock{\ell_{k-1}}{\ell_k}: mn\in S_k\}
  \end{equation}
  satisfies
  \begin{equation}\label{eq:retention}
    Q_{\ell_{k-1},\ell_k}\sum_{n\in F_k(m)}\frac1n
       \geq 1-2^{-k-2}.
  \end{equation}
\end{enumerate}
\end{definition}

The definition of an admissible family is independent of $A$.
Each condition at a fixed level involves finite sets, prime factorization, and finite reciprocal sums.
The multiplier $1$ is permitted in~\eqref{eq:fiber}.
The fibers may depend on the parent $m$.

\begin{lemma}\label{lem:family-properties}
Every admissible family satisfies $S_{k-1}\subseteq S_k$ for $k\geq1$.
Each fiber $F_k(m)$ is nonempty and divisor closed.
If $m',m\in S_{k-1}$ and $m'\mid m$, then $F_k(m)\subseteq F_k(m')$.
\end{lemma}

\begin{proof}
The right-hand side of~\eqref{eq:retention} is positive, so $F_k(m)$ is nonempty.
If $n\in F_k(m)$ and $d\mid n$, then $md\mid mn$.
Divisor closure of $S_k$ gives $md\in S_k$, and $d$ has the required prime support.
Thus $d\in F_k(m)$.
In particular $1\in F_k(m)$, so $m\in S_k$.
Finally, if $m'\mid m$ and $n\in F_k(m)$, then $m'n\mid mn$, giving $n\in F_k(m')$.
\end{proof}

\section{Prime valuations and the density transfer}

Let $V_1,V_2,\ldots$ be independent random variables with
\begin{equation}\label{eq:geometric}
  \Prob(V_i=e)=\left(1-\frac1{p_i}\right)p_i^{-e},
  \qquad e\in\Nzero.
\end{equation}
We work on the product probability space $\Omega=\Nzero^{\N}$ with these coordinate laws.
For $a\in\N$, write $a\mid V$ when $v_{p_i}(a)\leq V_i$ for every $i$.
This event depends on finitely many coordinates.
Define
\begin{equation}\label{eq:avoidance}
  C=C_A=\{V: a\nmid V\text{ for every }a\in A\}.
\end{equation}
The event $C$ is measurable and is downward closed in the coordinatewise order on $\Omega$.

For $\ell\in\Nzero$ and $0\leq r\leq s$, let
\[
  X_\ell=\prod_{i\leq\ell}p_i^{V_i},
  \qquad
  Y_{r,s}=\prod_{r<i\leq s}p_i^{V_i}.
\]
Independence gives the exact distributions
\begin{equation}\label{eq:block-law}
  \Prob(X_\ell=m)=\frac{Q_{0,\ell}}{m}\quad(m\in\Hblock{0}{\ell}),
  \qquad
  \Prob(Y_{r,s}=n)=\frac{Q_{r,s}}{n}\quad(n\in\Hblock{r}{s}).
\end{equation}
In particular, every finite-prefix cylinder has positive probability.
We denote the second probability in~\eqref{eq:block-law} by $\nu_{r,s}(n)$.
Then
\begin{equation}\label{eq:block-total}
  \sum_{n\in\Hblock{r}{s}}\nu_{r,s}(n)=1.
\end{equation}

\Needspace{8\baselineskip}
\begin{lemma}\label{lem:finite-density}
For every finite $F\subseteq\N$,
\begin{equation}\label{eq:finite-density}
  \Prob(a\mid V\text{ for some }a\in F)=d(M_F).
\end{equation}
\end{lemma}

\begin{proof}
For $e\in\Nzero$, summing~\eqref{eq:geometric} gives $\Prob(V_i\geq e)=p_i^{-e}$.
Consequently, for every nonempty finite $J\subseteq\N$,
\[
  \Prob(a\mid V\text{ for every }a\in J)
   =\prod_i p_i^{-v_{p_i}(\lcm(J))}
   =\frac1{\lcm(J)}.
\]
The positive integers divisible by every member of $J$ are precisely the multiples of $\lcm(J)$, and their density is the same quantity.
Applying finite inclusion-exclusion to both sides gives
\[
  \Prob(a\mid V\text{ for some }a\in F)
  =\sum_{\varnothing\neq J\subseteq F}
       \frac{(-1)^{|J|+1}}{\lcm(J)}
  =d(M_F).
\]
The assertion for $F=\varnothing$ is immediate.
\end{proof}

We use the classical finite-approximation theorem of Davenport and Erd\H{o}s~\cite{DavenportErdos1936}.

\begin{theorem}[Davenport--Erd\H{o}s]\label{thm:DE}
For every $A\subseteq\N$, with $A_t=A\cap\{1,\ldots,t\}$,
\begin{equation}\label{eq:DE}
  \underline d(M_A)=\lim_{t\to\infty}d(M_{A_t}).
\end{equation}
\end{theorem}

\begin{corollary}\label{cor:transfer}
For every $A\subseteq\N$,
\begin{equation}\label{eq:transfer}
  \Prob(C_A)=1-\underline d(M_A).
\end{equation}
In particular,
\begin{equation}\label{eq:zero-one-transfer}
  d(M_A)=1\quad\Longleftrightarrow\quad\Prob(C_A)=0.
\end{equation}
\end{corollary}

\begin{proof}
The events $\{V:a\mid V\text{ for some }a\in A_t\}$ increase to $C_A^c$.
Continuity from below, Lemma~\ref{lem:finite-density}, and Theorem~\ref{thm:DE} give~\eqref{eq:transfer}.
A subset of $\N$ with lower natural density one has natural density one, because its upper natural density is at most one.
This proves~\eqref{eq:zero-one-transfer}.
\end{proof}

\section{Positive avoidance from an admissible family}

\begin{proposition}\label{prop:positive-avoidance}
Suppose $(\ell_k,S_k)_{k\geq0}$ is admissible and $A\cap S_k=\varnothing$ for every $k$.
Then
\begin{equation}\label{eq:positive-bound}
  \Prob(C_A)\geq\frac34 Q_{0,\ell_0}>0.
\end{equation}
\end{proposition}

\begin{proof}
Set $E_k=\{V:X_{\ell_k}\in S_k\}$.
Since $S_0=\{1\}$, we have $\Prob(E_0)=Q_{0,\ell_0}$.
The projection condition in Definition~\ref{def:admissible} gives $E_k\subseteq E_{k-1}$.
For each $m\in S_{k-1}$, the new block $Y_{\ell_{k-1},\ell_k}$ is independent of $X_{\ell_{k-1}}$.
Equations~\eqref{eq:block-law} and~\eqref{eq:retention} therefore give
\begin{align}
  \Prob(E_k)
    &=\sum_{m\in S_{k-1}}\Prob(X_{\ell_{k-1}}=m)
       \sum_{n\in F_k(m)}\nu_{\ell_{k-1},\ell_k}(n)\notag\\
    &\geq (1-2^{-k-2})\Prob(E_{k-1}).\label{eq:survival-step}
\end{align}
For $u_j\in[0,1]$, induction gives
$\prod_{j=1}^k(1-u_j)\geq1-\sum_{j=1}^k u_j$.
Hence
\[
  \Prob(E_k)\geq Q_{0,\ell_0}\prod_{j=1}^k(1-2^{-j-2})
       \geq \frac34 Q_{0,\ell_0}.
\]
Continuity from above yields
\begin{equation}\label{eq:intersection-mass}
  \Prob\left(\bigcap_{k\geq0}E_k\right)\geq\frac34 Q_{0,\ell_0}.
\end{equation}

We claim that $\bigcap_{k\geq0}E_k\subseteq C_A$.
Indeed, let $V$ belong to this intersection and suppose $a\in A$ divides $V$.
The integer $a$ has finitely many prime divisors, and $\ell_k\to\infty$.
Thus $a\mid X_{\ell_k}$ for all sufficiently large $k$.
Since $X_{\ell_k}\in S_k$ and $S_k$ is divisor closed, it follows that $a\in S_k$.
This contradicts $A\cap S_k=\varnothing$.
The claim and~\eqref{eq:intersection-mass} prove~\eqref{eq:positive-bound}.
\end{proof}

\section{Finite-prefix extraction}

For $m\in\Hblock{0}{\ell}$, define
\begin{equation}\label{eq:conditional-survival}
  h_\ell(m)=\Prob(C_A\mid X_\ell=m).
\end{equation}
The conditioning is legitimate by~\eqref{eq:block-law}.
Equivalently, fix the first $\ell$ prime exponents to those of $m$ and sample the remaining coordinates according to~\eqref{eq:geometric}.
Let
\begin{equation}\label{eq:prefix-safe}
  C_\ell=\{V: a\nmid X_\ell\text{ for every }a\in A\}.
\end{equation}
Then $C_\ell\downarrow C_A$.
Indeed, any integer dividing $V$ divides $X_\ell$ once the prefix contains all of its prime divisors.

\begin{lemma}\label{lem:survival-properties}
For every $\ell\in\Nzero$,
\begin{enumerate}[label=\textup{(\roman*)},leftmargin=2.4em,itemsep=3pt]
  \item $0\leq h_\ell(m)\leq1$ for $m\in\Hblock{0}{\ell}$.
  \item If $m',m\in\Hblock{0}{\ell}$ and $m'\mid m$, then $h_\ell(m')\geq h_\ell(m)$.
  \item If some $a\in A$ divides $m$, then $h_\ell(m)=0$.
  \item $\Exp[h_\ell(X_\ell)]=\Prob(C_A)$.
\end{enumerate}
\end{lemma}

\begin{proof}
The first assertion follows from the definition as a conditional probability.
For the second, use the same future prime exponents with prefixes $m'$ and $m$.
The first complete valuation vector is coordinatewise no larger than the second.
Downward closure of $C_A$ shows that survival for the second implies survival for the first.
For the third, every completion of the prefix $m$ remains divisible by $a$.
The last assertion is the law of total probability applied to the countable partition $\{X_\ell=m\}$.
\end{proof}

\begin{lemma}\label{lem:zero-prefix}
If $c=\Prob(C_A)>0$, then $h_\ell(1)\to1$ as $\ell\to\infty$.
\end{lemma}

\begin{proof}
By Lemma~\ref{lem:survival-properties}, the function $h_\ell$ is at most $h_\ell(1)$ and vanishes on prefixes outside $C_\ell$.
Consequently,
\[
  c=\Exp[h_\ell(X_\ell)]\leq h_\ell(1)\Prob(C_\ell).
\]
Since $\Prob(C_\ell)\downarrow c>0$,
\[
  1\geq h_\ell(1)\geq\frac{c}{\Prob(C_\ell)}\longrightarrow1.\qedhere
\]
\end{proof}

\begin{lemma}\label{lem:finite-pruning}
Let $r\in\Nzero$, let $G\subseteq\Hblock{0}{r}$ be finite, and let $0<\delta<1$.
There exist $\ell>r$ and, for each $m\in G$, a finite set $N_m\subseteq\Hblock{r}{\ell}$ such that
\begin{equation}\label{eq:pruning-score}
  h_\ell(mn)\geq1-\delta\qquad(n\in N_m)
\end{equation}
and
\begin{equation}\label{eq:pruning-mass}
  \sum_{n\in N_m}\nu_{r,\ell}(n)>h_r(m)-2\delta.
\end{equation}
The cutoff $\ell$ is common to all $m\in G$.
\end{lemma}

\begin{proof}
Fix $m\in G$ and write $\Prob_m$ and $\Exp_m$ for probability and expectation conditional on $X_r=m$.
For $\ell\geq r$, we have $X_r=\pi_r(X_\ell)$.
The tower property and Lemma~\ref{lem:survival-properties}\textup{(iii)} give
\begin{align}
  D_\ell(m)
    &:=\Exp_m\bigl[(1-h_\ell(X_\ell))\one_{C_\ell}\bigr]\notag\\
    &=\Prob_m(C_\ell)-h_r(m).\label{eq:defect}
\end{align}
Indeed, $h_\ell(X_\ell)$ vanishes off $C_\ell$ and has conditional expectation $h_r(m)$ under $\Prob_m$.
Continuity from above gives
$\Prob_m(C_\ell)\downarrow\Prob_m(C_A)=h_r(m)$, and therefore $D_\ell(m)\to0$.
Since $G$ is finite, choose a common $\ell>r$ such that
\begin{equation}\label{eq:small-defect}
  D_\ell(m)<\delta^2\qquad(m\in G).
\end{equation}

On $C_\ell\cap\{h_\ell(X_\ell)<1-\delta\}$, the integrand in~\eqref{eq:defect} is greater than $\delta$.
It follows that
\[
  \Prob_m\bigl(C_\ell\cap\{h_\ell(X_\ell)<1-\delta\}\bigr)<\delta.
\]
The event $\{h_\ell(X_\ell)\geq1-\delta\}$ is contained in $C_\ell$, because $\delta<1$ and $h_\ell$ vanishes on unsafe prefixes.
Thus
\begin{equation}\label{eq:good-mass}
  \Prob_m\bigl(h_\ell(X_\ell)\geq1-\delta\bigr)
     >\Prob_m(C_\ell)-\delta
     \geq h_r(m)-\delta.
\end{equation}
Under $\Prob_m$, the new multiplier $Y_{r,\ell}$ has the law $\nu_{r,\ell}$, and $X_\ell=mY_{r,\ell}$.
Consequently, the probability on the left of~\eqref{eq:good-mass} is the sum of $\nu_{r,\ell}(n)$ over
\[
  \{n\in\Hblock{r}{\ell}:h_\ell(mn)\geq1-\delta\}.
\]
This is a convergent series of nonnegative terms.
Choose a finite subset $N_m$ whose sum is greater than that series minus $\delta$.
Equations~\eqref{eq:pruning-score} and~\eqref{eq:pruning-mass} follow.
\end{proof}

\section{Construction of an avoiding family}

\begin{proposition}\label{prop:construction}
If $\Prob(C_A)>0$, then there exists an admissible family $(\ell_k,S_k)_{k\geq0}$ such that $A\cap S_k=\varnothing$ for every $k$.
\end{proposition}

\begin{proof}
Set
\begin{equation}\label{eq:delta}
  \delta_k=2^{-k-5}\qquad(k\geq0).
\end{equation}
By Lemma~\ref{lem:zero-prefix}, choose $\ell_0$ with $h_{\ell_0}(1)\geq1-\delta_0$, and put $S_0=\{1\}$.
Lemma~\ref{lem:survival-properties}\textup{(iii)} gives $A\cap S_0=\varnothing$.
We construct the family inductively, maintaining
\begin{equation}\label{eq:inductive-score}
  h_{\ell_k}(s)\geq1-\delta_k\qquad(s\in S_k).
\end{equation}

Suppose that $\ell_{k-1}$ and the finite divisor-closed set $S_{k-1}\subseteq\Hblock{\ell_0}{\ell_{k-1}}$ have been chosen.
Apply Lemma~\ref{lem:finite-pruning} with
\[
  r=\ell_{k-1},\qquad G=S_{k-1},\qquad\delta=\delta_k.
\]
It gives $\ell_k>\ell_{k-1}$ and finite sets $N_m\subseteq\Hblock{\ell_{k-1}}{\ell_k}$ for $m\in S_{k-1}$.
Let
\begin{equation}\label{eq:construction}
  T_k=\bigcup_{m\in S_{k-1}}\{mn:n\in N_m\},
  \qquad
  S_k=\Div(T_k)=\{d\in\N:d\mid t\text{ for some }t\in T_k\}.
\end{equation}
The set $S_k$ is finite and divisor closed, and its elements use only primes in $(\ell_0,\ell_k]$.

If $d\mid mn$ with $m\in S_{k-1}$ and $n\in N_m$, then
$\pi_{\ell_{k-1}}(d)\mid m$, because $n$ uses only primes with indices greater than $\ell_{k-1}$.
Divisor closure of $S_{k-1}$ therefore gives
\[
  \pi_{\ell_{k-1}}(d)\in S_{k-1}.
\]
This verifies the projection condition.
Moreover, Lemma~\ref{lem:survival-properties}\textup{(ii)} and~\eqref{eq:pruning-score} imply
\[
  h_{\ell_k}(d)\geq h_{\ell_k}(mn)\geq1-\delta_k.
\]
Thus~\eqref{eq:inductive-score} holds at level $k$.
Since $1-\delta_k>0$, Lemma~\ref{lem:survival-properties}\textup{(iii)} shows that no element of $S_k$ has a divisor in $A$.
In particular $A\cap S_k=\varnothing$.

For every $m\in S_{k-1}$, the inclusion $mN_m\subseteq S_k$ gives $N_m\subseteq F_k(m)$.
Using~\eqref{eq:pruning-mass} and the induction hypothesis, we obtain
\begin{align}
  Q_{\ell_{k-1},\ell_k}\sum_{n\in F_k(m)}\frac1n
    &\geq\sum_{n\in N_m}\nu_{\ell_{k-1},\ell_k}(n)\notag\\
    &>h_{\ell_{k-1}}(m)-2\delta_k\notag\\
    &\geq1-\delta_{k-1}-2\delta_k
      =1-2^{-k-3}
      \geq1-2^{-k-2}.\label{eq:budget}
\end{align}
The lower bound is positive, so every $N_m$ is nonempty.
All conditions of Definition~\ref{def:admissible} have now been verified, completing the induction.
\end{proof}

\begin{proof}[Proof of Theorem~\ref{thm:main}]
Suppose $d(M_A)=1$.
By Corollary~\ref{cor:transfer}, $\Prob(C_A)=0$.
If an admissible family were disjoint from $A$ at every level, Proposition~\ref{prop:positive-avoidance} would give $\Prob(C_A)>0$.
Hence every admissible family meets $A$.

Conversely, suppose $A$ meets every admissible family.
If $M_A$ failed to have natural density one, Corollary~\ref{cor:transfer} would give $\Prob(C_A)>0$.
Proposition~\ref{prop:construction} would then produce an admissible family disjoint from $A$ at every level, a contradiction.
Therefore $d(M_A)=1$.
\end{proof}

\section*{Acknowledgments}
Artificial-intelligence tools provided by OpenAI assisted with mathematical exploration, checking arguments, and preparation of the manuscript.

\begin{thebibliography}{9}

\bibitem{DavenportErdos1936}
H.~Davenport and P.~Erd\H{o}s,
\emph{On sequences of positive integers},
Acta Arithmetica \textbf{2} (1936), 147--151.

\bibitem{Erdos1979}
P.~Erd\H{o}s,
\emph{Some unconventional problems in number theory},
Ast\'erisque \textbf{61} (1979), 73--82.

\end{thebibliography}
\end{document}
